\section{Resultados e Discussão}

Apresentar os resultados da pesquisa, podendo utilizar ilustrações como tabelas, figuras, quadros, fluxogramas e imagens.

\textbf{Normas para utilização das ilustrações:}

\begin{itemize}

\item Elas devem possuir títulos (cabeçalhos) localizados na parte superior, antecedidos da palavra que as designa (tabela, figura, esquema, fluxograma, imagem etc.), seguidos do número de ordem de ocorrência no texto, em algarismos arábicos, e travessão para separação do título.

\item As fontes consultadas são obrigatórias, mesmo que sejam de produção do próprio autor, devendo estar localizadas na parte inferior, contendo notas e outras informações necessárias à sua compreensão, quando aplicável.

\item As ilustrações devem ser centralizadas na página.

\item Os títulos e as fontes devem ficar alinhados à borda esquerda da ilustração. Caso ultrapassem sua largura, devem ser centralizados.

\item Deve-se utilizar fonte Arial, tamanho 12 para o título e tamanho 10 para a fonte, ambos sem ponto final.

\item Utilizar o termo \textbf{Figura} para desenhos e \textbf{Imagem} para fotografias.

\item As tabelas devem apresentar dados numéricos (valores, porcentagens, estatísticas etc.), enquanto os quadros devem apresentar dados qualitativos (ideias, descrições, resumos), sem uso de abreviações.

\item Caso a tabela não caiba em uma página, deve-se indicar ``continua'' ao final da primeira página e ``continuação'' no início da página seguinte, repetindo o título e o cabeçalho.

\end{itemize}

\vspace{0.5cm}

\textbf{Exemplo de Figura:}

% ================= FIGURA =================
\begin{figure}[H]
\caption{Figura 1 -- Exemplo de figura}
\centering
\includegraphics[width=0.5\textwidth]{Imagens/exemplo.jpg}

\raggedright
{\fontsize{10}{12}\selectfont Fonte: Elaborado pelo autor}
\end{figure}

\vspace{0.5cm}

\textbf{Exemplo de Tabela:}

% ================= TABELA =================
\begin{table}[H]
\caption{Tabela 1 -- Exemplo de tabela}
\centering
\begin{tabular}{lcc}
\hline
Item & Quantidade & Percentual \\
\hline
Teoria social & 22 & 7,9\% \\
Método & 34 & 12,3\% \\
Questão & 54 & 19,5\% \\
Raciocínio & 124 & 44,8\% \\
\hline
\end{tabular}

\raggedright
{\fontsize{10}{12}\selectfont Fonte: Adaptado de \cite{onu2023}}
\end{table}